\section*{Ключевые фрагменты кода}

\subsection*{Создание каналов и настройка наследования}
\begin{lstlisting}[language=C]
if (create_pipe_pair(&pipe1_read, &pipe1_write) != 0) handle_error("pipe1");
if (create_pipe_pair(&pipe2_read, &pipe2_write) != 0) handle_error("pipe2");

#ifdef _WIN32
SetHandleInformation(pipe1_write, HANDLE_FLAG_INHERIT, 0);
SetHandleInformation(pipe2_read, HANDLE_FLAG_INHERIT, 0);
#endif

process_info_t child = spawn_child("child.exe", filename, pipe1_read, pipe2_write);
\end{lstlisting}

\subsection*{Перенаправление потоков в Windows}
\begin{lstlisting}[language=C]
si.hStdInput = child_stdin;   // pipe1_read
si.hStdOutput = child_stdout; // pipe2_write
si.dwFlags |= STARTF_USESTDHANDLES;
CreateProcess(NULL, cmdline, NULL, NULL, TRUE, 0, NULL, NULL, &si, &pi);
\end{lstlisting}

\subsection*{Подсчет суммы в дочернем процессе}
\begin{lstlisting}[language=C]
while (scanf("%f", &value) == 1) {
    sum += value;
}
fprintf(file, "Sum: %.2f\n", sum);
printf("%.2f\n", sum);
\end{lstlisting}